\documentclass{article}

% Language setting
% Replace `english' with e.g. `spanish' to change the document language
\usepackage[english]{babel}

% Set page size and margins
% Replace `letterpaper' with `a4paper' for UK/EU standard size
\usepackage[letterpaper,top=2cm,bottom=2cm,left=3cm,right=3cm,marginparwidth=1.75cm]{geometry}
\usepackage[natbib=true, backend=biber, style=authoryear]{biblatex}
\addbibresource{references.bib}

% Useful packages

\usepackage{amsmath}
\usepackage{amsfonts}
\usepackage{graphicx}
\usepackage[colorlinks=true, allcolors=blue]{hyperref}
\usepackage{bm}
\usepackage{csquotes}



\newcommand{\cf}{cf.\xspace}
\newcommand{\eg}{e.g.\xspace}
\newcommand{\ie}{i.e.\xspace}

\title{WASP Software Engineering Assignment}
\author{
    Weijia Shao \\
    Chalmers University of Technology and University of Gothenburg
}

\date{August 2026}

\usepackage{minted}


\begin{document}
\maketitle

\section{Introduction}
Modern AI models are increasingly trained using distributed systems in which
multiple GPUs across several physical nodes cooperate on a single training
job. Frameworks such as PyTorch Distributed and communication libraries such
as NCCL coordinate computation and gradient exchange across these workers.
The underlying high-performance computing clusters are often multi-tenant:
different users may share a physical node while a scheduler such as SLURM
isolates their processes and allocated resources. However, this software-level
isolation does not necessarily partition shared hardware components such as
memory controllers, inter-socket links, PCIe paths, and network interfaces.

My research investigates the security and reliability consequences of this
shared infrastructure. In particular, I ask whether an unprivileged co-tenant
can deliberately degrade another user's distributed AI training job without
exploiting a software vulnerability or obtaining access to its processes,
models, or GPUs.


\section{Lecture principles}
\subsection{Principle 1: System Testing}
As discussed in the lecture on quality assurance, software testing means
systematically experimenting with an executable artifact to evaluate its
quality. Testing is therefore not limited to checking individual functions,
but can also evaluate integrated and complete systems under realistic
operating conditions. This is especially important when software, hardware,
requirements, and the surrounding environment interact in complex ways.

My research can be viewed as system-level stress testing of distributed AI
training infrastructure. Instead of testing only the model, the PyTorch code,
or the communication library, I execute the complete training system under an
adversarial but permitted condition: an unprivileged co-tenant aggressively
using shared host-memory resources. The relevant quality properties are
performance isolation, reliability, and availability rather than model
accuracy. A training job may remain functionally correct while still failing
as an engineered system if another tenant can increase its step time. My work
therefore extends AI testing beyond model outputs to the infrastructure
assumptions on which distributed training depends.

\subsection{Principle 2: Confirmation Bias and Disconfirming Evidence}
The lecture on behavioural software engineering highlighted confirmation bias:
engineers and researchers may favour evidence that supports their existing
explanations while overlooking contradictory observations. It also emphasized
traceability, generating alternative explanations, and explicitly searching
for disconfirming evidence as mitigation strategies.

This principle is important in my research because performance attribution in
distributed GPU systems is difficult. CPU execution, host memory, GPU kernels,
PCIe transfers, inter-socket communication, and network collectives proceed
concurrently, so an observed slowdown can have several plausible causes. After
seeing a longer AllReduce stage, for example, it would be easy to conclude
immediately that the network is the bottleneck while overlooking CPU
interference, measurement overhead, workload imbalance, or ordinary cluster
variability.

I therefore need to treat each proposed mechanism as a hypothesis that should
survive attempts to disprove it. This includes comparing memory and UPI
attacks, varying attacker placement and intensity, adding low-intensity and
baseline controls, testing workloads with different host-memory and
communication profiles, and comparing end-to-end timing with stage-level
measurements. Hardware performance counters and communication-library logs
provide additional evidence independent of the application timing. The goal is
not merely to show that training becomes slower, but to establish a credible
causal chain from contention in a particular shared hardware component to the
global slowdown of the distributed training job.

\section{Guest-lecture principles}
\subsection{Principle 1: Misuse Cases and Measurable Quality Requirements}
The guest lecture on requirements engineering emphasized that broad system
goals should be refined into concrete and measurable quality requirements. It
also introduced misuse cases, which describe behaviour that a system is
supposed to prevent.

This principle applies directly to my research. A general requirement such as
``the cluster shall provide secure and reliable multi-tenant AI training'' is
too vague to evaluate. It can instead be decomposed into measurable properties,
such as limiting the slowdown caused by a co-tenant, detecting abnormal
contention, and ensuring predictable performance isolation. My attack scenario
can also be formulated as a misuse case: a legitimate but adversarial tenant
uses only permitted CPU and memory operations to cause another tenant's
distributed training job to slow down substantially.

This formulation is important because the attacker does not bypass access
control or exploit a software vulnerability. Without an explicit misuse case,
the resulting slowdown may be treated merely as performance variability rather
than as a failure of security or reliability. The distinction between problem
space and solution space is also useful: my research first identifies the
missing system guarantee, while mechanisms such as throttling, monitoring,
scheduler changes, or hardware partitioning are possible solutions.

\subsection{Principle 2: Hidden Gaps in Responsibility and Control}
The guest lecture on human aspects of software engineering argued that
AI-mediated systems can expose hidden gaps in decisions, responsibilities, and
escalation paths. It also presented software engineering as more than code:
workflows, controls, roles, accountability, and organizational structures are
part of the engineered system.

My research exposes such a gap in multi-tenant AI infrastructure. A scheduler
may isolate processes and allocate CPUs and GPUs, but it does not necessarily
isolate shared memory controllers, inter-socket links, PCIe paths, or network
interfaces. Cluster operators may consider interference a workload issue,
framework developers may see it as an infrastructure issue, and users may
assume that resource allocation already implies sufficient isolation. As a
result, responsibility for preventing hardware-level noisy-neighbour attacks
is unclear.

By tracing the effect from a tenant's permitted memory accesses to hardware
contention, delayed communication, a straggling worker, and finally a global
training slowdown, my work makes this gap visible. The findings suggest that
mitigation requires not only technical changes, but also explicit monitoring,
performance guarantees, escalation procedures, and clearer responsibility
across cluster operators, schedulers, communication libraries, and AI users.

\section{Data scientists, software engineers, and AI engineers}

\subsection{Question 1}

The chapters characterise data scientists and software engineers as groups
with complementary but different backgrounds, workflows, and success
criteria. Data scientists tend to focus on data, statistical methods, model
design, experimentation, and predictive accuracy, often through exploratory
notebook-based workflows. Software engineers instead focus on delivering and
maintaining systems that meet user needs while balancing usability,
scalability, security, maintainability, time, and cost. I broadly agree with
this distinction, although it is an oversimplification rather than a strict
boundary. It reflects typical training and incentives, not a hierarchy of
skills, and many practitioners are better understood as T-shaped: specialised
in one area but able to collaborate across both.

\subsection{Question 2}

The system-wide framing applies directly to my research because the failure I
study is invisible from a model-centric perspective. The model architecture,
optimiser, and predictions may remain unchanged, while an unprivileged
co-tenant creates contention in shared memory controllers, inter-socket links,
or network paths. This delays one worker and consequently slows the complete
distributed training job. The vulnerability therefore emerges from the
interaction among the training algorithm, hardware, communication libraries,
and cluster scheduler. It demonstrates that performance, reliability, and
security are properties of the complete AI system, not of the model alone.

\subsection{Question 3}

I expect the roles to remain distinct specialisations, but not isolated skill
silos. Data scientists need enough knowledge of deployment, testing,
monitoring, and infrastructure to understand how their models behave in a
product. Software engineers need enough ML knowledge to reason about model
uncertainty, data quality, generalisation, and unavoidable prediction errors.
The likely outcome is interdisciplinary teams of T-shaped practitioners rather
than everyone becoming equally expert in both fields.

\subsection{Question 4}

AI Engineering is best understood as a partly distinct discipline at the
intersection of software engineering, MLOps, and ML systems work, rather than
as a rebranding of any one of them. It inherits software engineering concerns
such as requirements, architecture, quality assurance, operations, security,
and responsible design, while also addressing ML-specific issues such as data
dependence, statistical errors, expensive computation, and changing model
behaviour. As the chapters point out, MLOps focuses on making model pipelines
easier and more reliable to deploy, update, monitor, and operate, but this
pipeline view can remain model-centric if it stops at the interface between the
model service and the rest of the product. AI Engineering should therefore
take the broader system view advocated in the book, covering interactions
among models, software, hardware, users, and the operating environment. In my
research, this framing makes hardware-level interference an AI-engineering
problem rather than merely an HPC performance issue.

\subsection{Question 5}

Chapter 2 notes an emerging shift from building individual custom models
toward using foundation models and composing them through prompts, prompt
chains, and retrieval-augmented generation. Extending this observation,
foundation models, coding assistants, RAG systems, prompt pipelines, and
agentic workflows make the boundary between data science and software
engineering more permeable. Some work moves away from model training and
toward composing models, retrieval components, tools, and fallbacks into
reliable systems. This increases the importance of software-engineering skills
such as architecture, evaluation, monitoring, security, and failure handling,
while creating hybrid tasks such as prompt design, retrieval evaluation, and
agent orchestration. Coding assistants also allow data scientists to implement
software more easily and software engineers to prototype AI features, although
neither role can safely ignore the other's expertise. The distinction therefore
shifts from ``building models'' versus ``building software'' toward
collaboration across models, data, prompts, tools, infrastructure, and users.


\section{Paper analysis}
\subsection{Paper 1: Modeling Resilience of Collaborative AI Systems}
\paragraph{Core ideas and software-engineering importance.}

\cite{rimawi2024} propose a framework for modelling the resilience of
Collaborative Artificial Intelligence Systems (CAISs), where an AI component
works together with a human to achieve a shared goal. Their main idea is that
resilience should not be assessed only through a final accuracy value. Instead,
engineers should monitor how system performance evolves before, during, and
after a disruptive event. The framework divides execution into states,
including a steady state, a disruptive state, a recovered state, and a final
state after the source of disruption is removed. It then uses measurable
quantities to determine whether and how quickly the system returns to an
acceptable level of operation.

The paper demonstrates the framework using a collaborative robot that learns
online from human demonstrations to classify objects by colour. The robot
normally performs classifications autonomously, but falls back to human
assistance when its confidence is too low. The authors introduce the
Autonomous Classification Ratio as a measure of system performance and
deliberately disrupt the camera input by switching off its supporting lights.
They then track the resulting degradation and recovery. This is important to
software engineering because it treats resilience as an observable,
system-level quality property rather than as an informal expectation. It also
connects runtime monitoring with operational decisions: engineers can define
an acceptable threshold, detect when the system falls below it, and evaluate
the effort and time required for recovery.

\paragraph{Relation to my research.}

The paper relates closely to my research because both investigate how a local
disruptive event propagates into degradation of a larger AI system. In the
paper, the disruption changes sensor conditions and reduces the ability of an
online-learning robot to operate autonomously. In my work, an unprivileged
co-tenant creates contention in shared host hardware, delaying one worker in a
distributed training job and consequently slowing all workers.

An important mechanistic asymmetry sharpens this comparison. In the CAIS
setting the AI component itself adapts online and eventually closes the
disruptive state without external intervention. In synchronous distributed
training there is no equivalent self-healing: because the disruption sits in
the infrastructure rather than the model, the AllReduce barrier propagates
the damage across all workers instead of absorbing it, and recovery depends on
operator action rather than on continued learning.

My current evaluation focuses mainly on slowdown ratios, throughput loss, and
stage-level latency under attack. The resilience perspective suggests a
broader temporal analysis. For example, I could study how rapidly training
performance degrades after an attacker starts, whether the system reaches a
stable degraded state, and how quickly it recovers after the attacker stops.
This would distinguish resistance to disruption from recovery after disruption,
and would allow two systems with the same maximum slowdown to be compared
more meaningfully if one recovers much faster than the other.

\paragraph{Integration into a larger AI-intensive project.}

Consider a multi-tenant university platform used to train foundation models,
vision models, and graph neural networks on a shared GPU cluster. Its core
functionality includes distributed training, job scheduling, resource
allocation, communication through NCCL, and continuous collection of hardware
and application-level metrics. A resilience manager could apply the paper's
state-based model to each training job. Normal throughput would establish a
steady-state threshold; a sudden drop caused by contention would start a
disruptive state; and restoration of throughput after migration, throttling, or
attacker termination would indicate recovery.

My research would fit into this project by supplying both a validated
disruption mechanism and the hardware-counter signals needed to attribute it.
Hardware performance counters, GPU timing, communication logs, and topology
information could explain why the job crossed its resilience threshold. The
platform could then use this evidence to distinguish adversarial contention
from benign variation and trigger an appropriate response, such as job
relocation, resource throttling, or administrator investigation.

\paragraph{Adaptation of my research.}

To better support the paper's resilience idea, I would extend my experiments
from fixed attack intervals to complete disruption-and-recovery timelines.
Rather than reporting only baseline and attacked averages, I would measure the
onset time, degradation depth, accumulated lost training progress, recovery
time, and any residual performance effect after the attack ends. I would also
investigate whether different training frameworks or scheduling policies
degrade gracefully or abruptly. In the longer term, this could shift the
project from demonstrating an availability attack toward evaluating and
designing resilient distributed-training infrastructure.


\subsection{Paper 2: Trustworthy Distributed AI Systems: Robustness,
Privacy, and Governance}
\paragraph{Core ideas and software-engineering importance.}

\cite{wenqi2025} present a broad framework for trustworthy distributed AI based on
three main pillars: robustness, privacy, and fairness, supported by data and
model governance. The paper surveys threats that arise because distributed
learning involves multiple participants, communication channels, data sources,
and administrative domains. Its robustness taxonomy includes evasion attacks
during inference, poisoning and Byzantine attacks during training, and
irregular or non-IID data. Its privacy discussion covers attacks such as
gradient leakage and inference from model outputs, while its governance
discussion emphasizes monitoring, auditing, accountability, policy
enforcement, and compliance across the AI lifecycle.

The paper is important to software engineering because it treats
trustworthiness as an architectural and lifecycle concern rather than as a
property that can be added to a trained model afterwards. The authors argue
that trustworthy distributed AI should include built-in safeguards,
countermeasures, governance structures, and mechanisms for detecting,
preventing, and mitigating misuse. They also stress that robustness, privacy,
and fairness interact and therefore cannot be engineered entirely in
isolation.

\paragraph{Relation to my research.}

My research fits most directly under the paper's robustness and security
themes, but it also exposes a gap in its taxonomy. The survey mainly discusses
model- and data-level disruptions, including poisoning, Byzantine behaviour,
adversarial examples, privacy leakage, and irregular data distributions. My
work instead studies an infrastructure-level availability attack: a legitimate
co-tenant uses permitted memory operations to interfere with shared hardware
and slow distributed training without modifying the model, gradients, or
training data. Described in the survey's own vocabulary, this is a new row
in its threat table: attack target \emph{infrastructure / data-in-use}, attack
location \emph{co-tenant on the same physical node}, attack timing
\emph{training}, attack effect \emph{denial of service and cost inflation}.

A model can remain statistically correct and private while the system that
trains it becomes economically or operationally unusable. In expensive
distributed training, severe slowdown may increase cost, delay deployment,
waste allocated cluster time, or prevent time-critical model updates. This
motivates extending the survey's robustness pillar with an
infrastructure-integrity dimension alongside its existing algorithmic
categories.

\paragraph{Integration into a larger AI-intensive project.}

The same multi-tenant AI training platform can illustrate how the paper's ideas
would be applied. The platform could maintain a trustworthiness layer covering
model integrity, privacy, fairness, infrastructure availability, and
governance. Existing mechanisms might detect poisoned updates, protect
gradients, audit data use, and evaluate model fairness. My research would add
an infrastructure-security component that detects abnormal memory-controller,
inter-socket, PCIe, or network contention and attributes it to co-located
workloads.

The platform could combine these signals in a governance dashboard used by
cluster operators, security engineers, and AI researchers. Policies could
define acceptable interference, specify which tenants may share particular
hardware paths, and determine when a job should be migrated or an attacker
should be throttled. This would operationalize the paper's argument that
trustworthiness requires both technical safeguards and enforceable governance,
rather than relying only on voluntary good behaviour by participants.

\paragraph{Adaptation of my research.}

To better support the paper's vision, I would broaden my threat model from one
specific memory-contention attack to a taxonomy of infrastructure-level threats
against distributed AI. This could include contention on memory controllers,
inter-socket links, PCIe, storage, and network interfaces, as well as
accidental and adversarial noisy neighbours. I would then classify each threat
according to detection, prevention, mitigation, and recovery.

I would also empirically probe how the survey's own defences interact with
this new threat class. A concrete hypothesis worth testing is that
privacy-enhancing mechanisms such as TEE-hosted local training or homomorphic
encryption, which add memory-traffic and computation overhead, may in fact
amplify sensitivity to co-tenant contention rather than mitigate it. A
negative interaction result of this kind would be a novel contribution to the
trustworthy-AI literature.

Finally, I would connect technical measurements to governance requirements.
The project could define measurable service-level objectives for performance
isolation and develop evidence that a cluster operator could use for auditing
or policy enforcement. In the long term, this would position my work as part
of trustworthy AI infrastructure: not only showing that attacks exist, but
helping specify, monitor, and enforce the guarantees required for secure and
reliable distributed training.


\section{Evidence, originality, and GenAI-use transparency}
The main sources used were the course lecture slides, the two assigned
chapters of the CMU \emph{Machine Learning in Production} book, and the two
peer-reviewed papers reviewed in Section~5. The lecture slides were used to
identify and interpret the selected software-engineering principles. The book
chapters were used for the comparison of data scientists, software engineers,
and AI engineers. The CAIN paper is a peer-reviewed conference paper, while
the second source is a peer-reviewed survey in \emph{ACM Computing Surveys}.
I place the greatest trust in the peer-reviewed papers and book chapters
because their claims and terminology are presented with fuller context and
supporting evidence. ChatGPT was used to refine the language of the report. 



\printbibliography
\end{document}